% sec2_prelim.tex: Section 2 (Preliminaries).
% Labels defined here: sec:prelim, sec:prelim-op, sec:prelim-rates, sec:prelim-eq, sec:tool-wlxz, sec:tool-bessel,
%   sec:tool-cd, sec:tool-dlmf, eq:rec-monic, eq:rec-orth, eq:coef-int, def:regular, eq:regular-facts,
%   rem:majorant, prop:equilibrium, eq:Y-solution, eq:Y-extract, eq:CD-moments.

\section{Preliminaries}\label{sec:prelim}

This section fixes conventions and collects the inputs of the proofs. The only input from the companion
paper~\cite{Paper1} is the explicit equilibrium measure (\S\ref{sec:prelim-eq}). Every external result used later is
stated in \S\S\ref{sec:tool-wlxz}--\ref{sec:tool-dlmf}, with its hypotheses and the equation numbers of the source,
so that the later sections can refer to it by subsection and part.

\emph{Notation used throughout.} The data are $p>0$, $s>1$ and $\ell$; the function $\Lambda$, the modulus
$\varepsilon_\Lambda$ and the notion ``analytic of angle $\vartheta$'' are those of Definition~\ref{def:analytic}, and
$c=c(p)$ is given by \eqref{eq:cp}. The Bessel order of the hard edge is $\alpha:=(s-1-p)/p$, so that $\alpha>-1$ for
every $s>1$. Sections~\ref{sec:frame} and~\ref{sec:outer} hold for every $\alpha>-1$ and every $\ell$ that is analytic
of some angle; Sections~\ref{sec:hardedge} and~\ref{sec:extraction} assume $|\alpha|<1$, that is, $1<s<1+2p$; and
Sections~\ref{sec:classical} and~\ref{sec:second} assume $p>1$ and $\ell\equiv1$. The Riemann--Hilbert notation is
introduced in Section~\ref{sec:frame}, and Table~\ref{tab:geometry} lists the quantities that depend on whether $p>1$,
$p=1$ or $p<1$. Symbols used only inside one proof are local to it. We write $\mathbb D(z_0,r)$ for the open disc of
radius $r$ about $z_0$, and $\C_\pm$ for the open upper and lower half-planes.

\subsection{Orthogonal polynomials and recurrence coefficients}\label{sec:prelim-op}

Let $\sigma$ be a positive Borel measure on $[0,\infty)$ with infinite support and finite moments. Its Hankel
determinants are positive, so the monic orthogonal polynomials $P_n$ ($\deg P_n=n$) exist for every $n\ge0$. Put
$h_n:=\int P_n^2\diff\sigma>0$, $P_{-1}:=0$ and $b_0:=0$. The polynomials satisfy the three-term recurrence
\begin{equation}\label{eq:rec-monic}
  xP_n=P_{n+1}+a_nP_n+b_n^2P_{n-1}\qquad(n\ge0),\qquad b_n^2=\frac{h_n}{h_{n-1}}\quad(n\ge1),
\end{equation}
with $a_n\in\R$ and $b_n>0$ for $n\ge1$. We call $a_n$ and $b_n$ the \emph{recurrence coefficients} of
$\sigma$. Multiplying $\sigma$ by a positive constant does not change them. The orthonormal polynomials $p_n:=P_n/\sqrt{h_n}$
satisfy
\begin{equation}\label{eq:rec-orth}
  xp_n=b_{n+1}p_{n+1}+a_np_n+b_np_{n-1}\qquad(n\ge0),
\end{equation}
as one sees by dividing \eqref{eq:rec-monic} by $\sqrt{h_n}$ and using $b_n^2\sqrt{h_{n-1}/h_n}=b_n$. Since $xp_j$ has
degree $j+1$, orthonormality gives
\begin{equation}\label{eq:coef-int}
  a_n=\int xp_n^2\diff\sigma,\qquad b_n=\int xp_{n-1}p_n\diff\sigma,\qquad \int xp_np_j\diff\sigma=0\quad(j\le n-2).
\end{equation}
In matrix form, \eqref{eq:rec-orth} says that the Jacobi matrix of $\sigma$ is the symmetric tridiagonal matrix
$\mathcal J(\sigma)$ with diagonal $(a_0,a_1,\dots)$ and off-diagonal $(b_1,b_2,\dots)$. The monic form
\eqref{eq:rec-monic} corresponds to the non-symmetric matrix $J=J(\sigma)$ with $J_{n,n}=a_n$, $J_{n,n+1}=1$ and
$J_{n+1,n}=b_{n+1}^2$; this is the form used in Section~\ref{sec:shift}. The ratio $a_n/(2b_n)$ always pairs the
diagonal entry at $P_n$ with the off-diagonal entry $b_n$ between $P_{n-1}$ and $P_n$; both are read off from the same
Riemann--Hilbert matrix of degree $n$ (Proposition~\ref{prop:frame}(viii)).

If $\ell$ is analytic of angle $\vartheta$ for some $\vartheta>0$, then $\ell(k)\le C_\eta k^\eta$ for every $\eta>0$
by the Potter bound of Lemma~\ref{lem:pottersector}(i). Taking $\eta<s-1$ shows that $\sigma_{p,s,\ell}$ is a finite
measure on $(0,1]$ with infinitely many atoms, so everything above applies to it.

\subsection{Regular rates}\label{sec:prelim-rates}

\begin{definition}\label{def:regular}
A sequence $(\eta_n)_{n\ge1}$ of positive numbers is a \emph{regular rate} if $\eta_n$ is non-increasing, $n\eta_n$ is
non-decreasing, and $\eta_n\to0$.
\end{definition}

The basic example is $\eta_n=1/n$. Every regular rate satisfies
\begin{equation}\label{eq:regular-facts}
  \frac1n\le\frac{\eta_n}{\eta_1},\qquad \eta_{n-1}\le2\eta_n\quad(n\ge2),\qquad \sum_n\eta_n=\infty .
\end{equation}
The first holds because $n\eta_n\ge\eta_1$; the second because $(n-1)\eta_{n-1}\le n\eta_n$; the third follows from
the first.

\begin{remark}[Regular majorants]\label{rem:majorant}
Every non-increasing sequence $\mathsf v_n>0$ with $\mathsf v_n\to0$ is dominated by a regular rate, namely
$\eta_n:=\max_{m\le n}(m\mathsf v_m)/n$. Here $\eta_n\ge\mathsf v_n$ and $n\eta_n$ is non-decreasing by construction;
$\eta_{n+1}\le\max\bigl(\eta_n\,n/(n+1),\mathsf v_{n+1}\bigr)\le\eta_n$ because $\mathsf v_{n+1}\le\mathsf v_n\le\eta_n$;
and $\eta_n\to0$, since $m\mathsf v_m/n\le\mathsf v_1/\sqrt n$ for $m\le\sqrt n$ and $m\mathsf v_m/n\le\mathsf
v_{\lceil\sqrt n\rceil}$ for $\sqrt n<m\le n$. So Theorem~\ref{thm:shift} applies to any monotone rate. In this paper we
use it only with $\eta_n=1/n$.
\end{remark}

\subsection{The equilibrium measure}\label{sec:prelim-eq}

The Riemann--Hilbert analysis of Sections~\ref{sec:frame}--\ref{sec:second} is built on the equilibrium measure of the
constrained problem solved in~\cite{Paper1}. We use only its explicit form and the equality of its potential on the
band; the optimality statement \cite[Theorem~1.3(iii)]{Paper1} is not used in this paper.

Let $c=c(p)$, let $\gamma_p:=\log c(p)+p$, and let $\lambda_p(x):=\frac1px^{-1-1/p}$ be the density of the nodes
$k^{-p}$ (the image of Lebesgue measure on $(0,\infty)$ under $u\mapsto u^{-p}$). For $y\in[0,1]$ and $a',b'>0$ write
$\mathcal B_y(a',b'):=B(a',b')^{-1}\int_0^yr^{a'-1}(1-r)^{b'-1}\diff r$ for the regularised incomplete beta function,
where $B$ is Euler's Beta function. For $\beta>0$ let $\mathrm{arc}_\beta$ be the arcsine law on $[0,\beta]$, with
density $1/(\pi\sqrt{y(\beta-y)})$, and put $\beta_\tau:=4c\tau^{-p}$ for $0<\tau\le1$.

\begin{proposition}\label{prop:equilibrium}
Let $p>0$ and let $\mu^*:=\int_0^1\mathrm{arc}_{\beta_\tau}\diff\tau$, a probability measure on $(0,\infty)$.
\begin{enumerate}
\item[\textup{(i)}] $\mu^*$ has the density
  \[
    \psi^*(x)=\lambda_p(x)\,\mathcal B_{\min(1,\,x/4c)}\bigl(\tfrac1p+\tfrac12,\tfrac12\bigr),\qquad x>0 .
  \]
  In particular $\psi^*=\lambda_p$ on $[4c,\infty)$ (saturation) and $0<\psi^*<\lambda_p$ on $(0,4c)$.
\item[\textup{(ii)}] The function $U^*(x):=\int\log|x-y|\diff\mu^*(y)$ satisfies $U^*(x)=\gamma_p$ for $0<x\le4c$, and
  $U^*(x)\ge\gamma_p$ for all $x>0$.
\item[\textup{(iii)}] \textup{(}Hard edge.\textup{)} As $x\to0^+$,
  \[
    \psi^*(x)=C_*x^{-1/2}\bigl(1+O(x)\bigr),\qquad C_*:=\frac{(4c)^{-1/p-1/2}}{p\,a_*B(a_*,\frac12)},\quad
    a_*:=\frac1p+\frac12 .
  \]
\end{enumerate}
\end{proposition}

\begin{proof}
Parts (i) and (ii) are \cite[Theorem~1.3(i)--(ii)]{Paper1}, transported from the variable $u$ to $x=u^{-p}$. There,
$\mu^*_u$ is the image of $\mu^*$ under $x\mapsto x^{-1/p}$ and has density
$\rho^*(u)=\mathcal B_{\min(1,u^{-p}/4c)}(\frac1p+\frac12,\frac12)$, with $\rho^*=1$ exactly on $(0,(4c)^{-1/p}]$ and
$0<\rho^*<1$ on $((4c)^{-1/p},\infty)$. Since $|\diff u/\diff x|=\lambda_p(x)$, the density of $\mu^*$ in $x$ is
$\lambda_p(x)\rho^*(x^{-1/p})$, which is (i). Part (ii) is quoted verbatim, with $\gamma_p=\log c+p$. For (iii), let
$y=x/4c\le1$. Then $\int_0^yr^{a_*-1}(1-r)^{-1/2}\diff r=\frac{y^{a_*}}{a_*}(1+O(y))$, so
$\psi^*(x)=\frac1px^{-1-1/p}(x/4c)^{1/p+1/2}\bigl(a_*B(a_*,\frac12)\bigr)^{-1}(1+O(x))$, which is (iii).
\end{proof}

The point $x=4c$ is the soft edge, $(0,4c)$ is the band, and $x=0$ is the hard edge, where $\mu^*$ has the
inverse-square-root density of (iii).

\subsection{The discrete interpolation problem of Wu, Lin, Xu and Zhao}\label{sec:tool-wlxz}

This subsection and the next three state the external results used in Sections~\ref{sec:frame}--\ref{sec:second} and
the appendices. They are needed only from Section~\ref{sec:frame} on. Throughout, the Cauchy transform on an oriented
contour $\Sigma$ is
\[
  \Cauchy f(z):=\frac1{2\pi i}\int_\Sigma\frac{f(s)\diff s}{s-z},\qquad z\in\C\setminus\Sigma,
\]
and $\Cauchy_\pm f$ are its non-tangential boundary values from the left ($+$) and the right ($-$) of $\Sigma$, so that
$\Cauchy_+-\Cauchy_-=1$. For a jump matrix $v$ on $\Sigma$ we write $\Cauchy_vf:=\Cauchy_-\bigl(f(v-I)\bigr)$. The
matrices $\sigma_1,\sigma_2,\sigma_3$ are the Pauli matrices, and $\sigma_3=\mathrm{diag}(1,-1)$.

Wu, Lin, Xu and Zhao~\cite{WLXZ} develop the Riemann--Hilbert method for discrete orthogonal polynomials on infinitely
many nodes with an accumulation point, for the Tricomi--Carlitz polynomials. We apply none of their theorems to our
measures; we use the formulation of the problem and the order of the transformations, and our versions are written out
in Sections~\ref{sec:frame} and~\ref{sec:outer}.

\emph{(a) The problem and its solution.} Let $\sigma=\sum_{k\ge1}w_k\delta_{x_k}$ with distinct nodes
$x_k\in(0,1]$, $w_k>0$ and $\sum_kw_k<\infty$, and fix $n\ge1$. Following \cite[(2.14) and Theorem~2.3]{WLXZ}, put
\begin{equation}\label{eq:Y-solution}
  Y(z):=\begin{pmatrix}P_n(z) & \displaystyle\sum_k\frac{w_kP_n(x_k)}{z-x_k}\\[3mm]
  h_{n-1}^{-1}P_{n-1}(z) & \displaystyle h_{n-1}^{-1}\sum_k\frac{w_kP_{n-1}(x_k)}{z-x_k}\end{pmatrix}.
\end{equation}
Then $Y$ is analytic in $\C\setminus(\{x_k\}_k\cup\{0\})$; at each node it has a simple pole with
$\Res_{z=x_k}Y(z)=\lim_{z\to x_k}Y(z)\bigl(\begin{smallmatrix}0&w_k\\0&0\end{smallmatrix}\bigr)$ (the residue
condition \cite[(2.14)]{WLXZ}); and $Y(z)z^{-n\sigma_3}=I+Y_1/z+Y_2/z^2+O(z^{-3})$ as $z\to\infty$, with
\begin{equation}\label{eq:Y-extract}
  (Y_1)_{12}(Y_1)_{21}=b_n^2,\qquad \frac{(Y_2)_{12}}{(Y_1)_{12}}-(Y_1)_{22}=a_n .
\end{equation}
\emph{Proof.} The sums converge locally uniformly off $\{x_k\}\cup\{0\}$, which gives analyticity, the simple poles and
$\Res_{x_k}Y_{j2}=w_kY_{j1}(x_k)$. Write $P_m(x)=x^m+\pi_mx^{m-1}+\dots$. For $|z|>2$, expanding $(z-x)^{-1}$ gives
$\sum_kw_kP_m(x_k)/(z-x_k)=\sum_{i\ge0}z^{-i-1}\int P_mx^i\diff\sigma$, and $\int P_mx^i\diff\sigma$ equals $0$ for $i<m$,
$h_m$ for $i=m$ and $-\pi_{m+1}h_m$ for $i=m+1$. Hence $(Y_1)_{12}=h_n$, $(Y_1)_{21}=h_{n-1}^{-1}$,
$(Y_2)_{12}=-\pi_{n+1}h_n$ and $(Y_1)_{22}=-\pi_n$; and $\pi_n=\pi_{n+1}+a_n$ by \eqref{eq:rec-monic}. \qed

\emph{(b) Transformations.} The transformations of Section~\ref{sec:frame} follow \cite[\S\S3--5]{WLXZ}: removal of the
poles by triangular factors built from $\Pi=\sin\theta/\theta'$ and $e^{\pm i\theta}$ \cite[(3.9)--(3.15)]{WLXZ},
normalisation by the $g$-function \cite[(4.1)]{WLXZ}, the lens and saturated-region factors $V_\pm$, $V^\triangle_\pm$
\cite[(4.12)--(4.13)]{WLXZ}, and conjugation by an outer scalar \cite[(5.1)]{WLXZ}. In \cite{WLXZ} these are written in
a symmetric variable with nodes $\pm x_k$; we use only the pattern, and every jump is recomputed in
Section~\ref{sec:frame} for our one-sided problem.

\emph{(c) The soft edge.} Our Airy parametrix (Lemma~\ref{lem:airy}) follows the pattern of \cite[\S6.1, (6.4),
(6.10)--(6.12)]{WLXZ}; it is constructed from the Airy function directly (\S\ref{sec:tool-dlmf}), and no statement of
\cite{WLXZ} is used.

\subsection{The Bessel model of Kuijlaars, McLaughlin, Van Assche and Vanlessen}\label{sec:tool-bessel}

We use the explicit Bessel model of \cite[\S6]{KMVV} and its algebraic properties. In Sections~\ref{sec:hardedge}
and~\ref{sec:extraction} we apply none of their theorems on orthogonal polynomials, so their standing hypothesis on the
weight \cite[(1.1)]{KMVV} plays no role there. Section~\ref{sec:classical} uses the construction of part~(c) at the
left endpoint; the hypotheses it needs are checked there.

\emph{(a) The model.} Let $\alpha>-1$ and let $\Sigma_\Psi$ consist of the rays $\gamma_1:\arg\zeta=2\pi/3$,
$\gamma_2:\arg\zeta=\pi$ and $\gamma_3:\arg\zeta=-2\pi/3$, all oriented towards the origin, with the $+$ side on the
left \cite[Figure~3]{KMVV}. The function
\[
  \Psi_\alpha(\zeta)=\begin{pmatrix}I_\alpha(2\zeta^{1/2}) & \frac i\pi K_\alpha(2\zeta^{1/2})\\[1mm]
  2\pi i\zeta^{1/2}I_\alpha'(2\zeta^{1/2}) & -2\zeta^{1/2}K_\alpha'(2\zeta^{1/2})\end{pmatrix},\qquad|\arg\zeta|<2\pi/3
\]
\cite[(6.23)]{KMVV}, continued to $2\pi/3<\pm\arg\zeta<\pi$ by the Hankel-function formulas \cite[(6.24)--(6.25)]{KMVV},
solves the Bessel model problem \cite[Theorem~6.3]{KMVV}, with jumps
\[
  \begin{pmatrix}1&0\\e^{\alpha\pi i}&1\end{pmatrix}\ \text{on }\gamma_1,\qquad
  \begin{pmatrix}0&1\\-1&0\end{pmatrix}\ \text{on }\gamma_2,\qquad
  \begin{pmatrix}1&0\\e^{-\alpha\pi i}&1\end{pmatrix}\ \text{on }\gamma_3
\]
\cite[(6.16)--(6.18)]{KMVV}. Uniformly as $\zeta\to\infty$ in $\C\setminus\Sigma_\Psi$,
\[
  \Psi_\alpha(\zeta)=(2\pi\zeta^{1/2})^{-\sigma_3/2}\frac1{\sqrt2}
  \begin{pmatrix}1+O(\zeta^{-1/2})&i+O(\zeta^{-1/2})\\ i+O(\zeta^{-1/2})&1+O(\zeta^{-1/2})\end{pmatrix}e^{2\zeta^{1/2}\sigma_3}
\]
\cite[(6.28)]{KMVV}, which follows from the large-argument expansions (\S\ref{sec:tool-dlmf}) of the Bessel and Hankel
functions in (6.23)--(6.25); this also covers the rays moved by (b) (\S\ref{sec:he-model}).

\emph{(b) Angle freedom.} By the remark after \cite[(6.21)]{KMVV}, the angle $2\pi/3$ may be replaced by any angle in
$(0,\pi)$; the lens contours in the disc are then the preimages of the rays under $\zeta=n^2f(z)$, and they do not
depend on $n$. We use rays at $\pm(\pi-\thlip)$.

\emph{(c) The left endpoint.} Our hard edge has the band to its right, as the endpoint $-1$ of \cite{KMVV}. There the
conformal map reverses orientation, and the reversal is absorbed by conjugation with $\sigma_3$:
$\tilde\Psi(\zeta)=\sigma_3\Psi(\zeta;\alpha\to\beta)\sigma_3$ \cite[(6.51)]{KMVV}, with prefactor
\[
  \tilde E_n(z)=N(z)\tilde W(z)^{\sigma_3}\tfrac1{\sqrt2}\begin{pmatrix}1&i\\i&1\end{pmatrix}
  \tilde f(z)^{\sigma_3/4}(2\pi n)^{\sigma_3/2}
\]
\cite[(6.53)]{KMVV}, where $N$, $\tilde W$ and $\tilde f$ are the global parametrix, an analytic square root of the
weight and the local conformal map of \cite[(6.42)--(6.50)]{KMVV} (compare the right-endpoint formulas
\cite[(6.27), (6.31)]{KMVV} and $\det E_n\equiv1$ \cite[(6.33)]{KMVV}). The prefactor is analytic by removability: its
entries have at most a square-root singularity at the endpoint and it is analytic in a punctured neighbourhood
\cite[proof of Proposition~6.5]{KMVV}. In Section~\ref{sec:hardedge} this argument is applied only to a prefactor that
does not involve the weight.

\subsection{The comparison argument of Conway and Deift}\label{sec:tool-cd}

The following are from \cite[\S4]{CD}, whose Cauchy operator is normalised as ours. Lebesgue exponents are written
$r,r'$ (\cite{CD} writes $p,q$).

\emph{(a) Boundedness and the class $\dC$} \cite[\S4.1]{CD}. $\Cauchy_\pm$ are bounded on $L^2(\Sigma)$ (and on
$L^r(\Sigma)$, $1<r<\infty$) if and only if $\Sigma$ is Carleson. A pair $A_\pm\in L^r(\Sigma)$ belongs to
$\dC(L^r(\Sigma))$ if $A_\pm=\Cauchy_\pm f$ for some (unique) $f\in L^r(\Sigma)$; we write $A_\pm\in I+\dC(L^r)$ if
$A_\pm-I\in\dC(L^r)$.

\emph{(b) Theorem~4.4 of \cite{CD}.} \emph{Let $v_A,v_B\in L^\infty(\Sigma)$ with $v_B^{-1}\in L^\infty(\Sigma)$. Let
$A_\pm\in I+\dC(L^r(\Sigma))$ satisfy $A_+=A_-v_A$ on $\Sigma$, and let $B_\pm\in I+\dC(L^r(\Sigma))$ satisfy
$B_+=B_-v_B$ on $\Sigma$, with $B_\pm^{-1}\in I+\dC(L^{r'}(\Sigma))$, $1/r+1/r'=1$. If $A,B$ are the analytic
extensions of $A_\pm,B_\pm$, then}
\[
  AB^{-1}=I+\Cauchy\bigl(A_-(v_Av_B^{-1}-I)B_-^{-1}\bigr).
\]
\emph{Consequence (moments at infinity).} If in addition $\Sigma$ is bounded and $A=I+A_1/z+A_2/z^2+O(z^{-3})$,
$B=I+B_1/z+B_2/z^2+O(z^{-3})$ as $z\to\infty$, then with $\mathcal E:=A_-(v_A-v_B)v_B^{-1}B_-^{-1}\in L^1(\Sigma)$
(take $r=r'=2$),
\begin{equation}\label{eq:CD-moments}
  A_1-B_1=-\frac1{2\pi i}\int_\Sigma\mathcal E(s)\diff s,\qquad
  A_2-B_2=-\frac1{2\pi i}\int_\Sigma s\mathcal E(s)\diff s+(A_1-B_1)B_1 .
\end{equation}
Indeed $v_Av_B^{-1}-I=(v_A-v_B)v_B^{-1}$; for $|z|>\max_\Sigma|s|$ expand $1/(s-z)=-\sum_{j\ge0}s^jz^{-j-1}$ in
$\Cauchy\mathcal E(z)$ and compare with $AB^{-1}=I+(A_1-B_1)z^{-1}+\bigl(A_2-B_2-(A_1-B_1)B_1\bigr)z^{-2}+O(z^{-3})$.
This is the mechanism of \cite[Proposition~4.9]{CD}, whose proof identifies $\mu=Q_-$ and $\tilde\mu=\tilde Q_-$.
Formula (4.14) of \cite{CD} is printed without the factor $v_B^{-1}$; it agrees with \eqref{eq:CD-moments} when
$(v_A-v_B)v_B^{-1}=v_A-v_B$, as in their setting. We use the form derived from Theorem~4.4.

\emph{(c) Solution formula} \cite[\S4.2]{CD}. If $1-\Cauchy_{v}$ is invertible on $L^2(\Sigma)$, then
$\mu:=(1-\Cauchy_v)^{-1}I$ gives the unique solution $X(z)=I+\Cauchy(\mu(v-I))(z)$ of the Riemann--Hilbert problem
normalised at infinity with jump $v$.

\emph{(d) The Neumann-series step} \cite[proof of Theorem~4.7, (4.12)--(4.13)]{CD}. If
$\|(1-\Cauchy_{\tilde v})^{-1}\|_{L^2\to L^2}\le K$ and $\|\Cauchy_v-\Cauchy_{\tilde v}\|_{L^2\to L^2}\le1/(2K)$, then
\[
  (1-\Cauchy_v)^{-1}=(1-\Cauchy_{\tilde v})^{-1}\bigl(1-(\Cauchy_v-\Cauchy_{\tilde v})(1-\Cauchy_{\tilde v})^{-1}\bigr)^{-1}
\]
exists and has norm at most $2K$, since the inner factor is $1$ minus an operator of norm at most $\frac12$. In
\cite{CD} the smallness comes from $\|v-\tilde v\|_{L^\infty(\Sigma)}\to0$.

\emph{(e) Proposition~4.6 of \cite{CD}} (attributed there to Deift and Zhou). The operators $(1-\Cauchy_{h_n})^{-1}$
are bounded uniformly in $n$ if and only if the inhomogeneous problem $m_+=m_-h_n+g$ on $\Sigma$,
$m_\pm\in\dC(L^2(\Sigma))$, is uniquely solvable for every $g\in L^2(\Sigma)$ with $\|m_\pm\|_{L^2}\le C\|g\|_{L^2}$, $C$
independent of $n$ and $g$. We use only the direction ``solvable with bounds $\Rightarrow$ bounded inverse'', and we
re-prove it in Appendix~\ref{app:resolvent}.

We do not use \cite[Theorem~4.5]{CD}, whose proof is specific to the Legendre comparison problem.

\subsection{Special functions}\label{sec:tool-dlmf}

All formulas are from \cite{DLMF}, with principal branches throughout. ``Uniformly in closed subsectors'' of a sector
means uniformly in every closed sector contained in it.

\emph{(a) The Gamma function.} (5.5.1) $\Gamma(z+1)=z\Gamma(z)$; (5.5.3) $\Gamma(z)\Gamma(1-z)=\pi/\sin(\pi z)$,
$z\notin\mathbb Z$; (5.11.1) $\mathrm{Ln}\,\Gamma(z)\sim(z-\frac12)\ln z-z+\frac12\ln(2\pi)+\sum_{k\ge1}
\frac{B_{2k}}{2k(2k-1)z^{2k-1}}$ as $z\to\infty$, uniformly in closed subsectors of $|\mathrm{ph}\,z|<\pi$, where by
\S5.11(ii) the remainder after the term $k=m-1$ is at most $\sec^{2m}(\frac12\mathrm{ph}\,z)$ times the first neglected
term; (5.12.3), (5.12.1) $\int_0^\infty t^{a'-1}(1+t)^{-a'-b'}\diff t=B(a',b')=\Gamma(a')\Gamma(b')/\Gamma(a'+b')$ for
$\re a',\re b'>0$.

\emph{(b) Airy functions} (\S\S9.2, 9.7). $\mathrm{Ai}$ is entire, and $\mathrm{Ai}(z)$, $\mathrm{Ai}(ze^{\mp2\pi
i/3})$ solve $w''=zw$; (9.2.8) $\mathcal W\{\mathrm{Ai}(z),\mathrm{Ai}(ze^{\mp2\pi i/3})\}=e^{\pm\pi i/6}/(2\pi)$;
(9.2.12) $\mathrm{Ai}(z)+e^{-2\pi i/3}\mathrm{Ai}(ze^{-2\pi i/3})+e^{2\pi i/3}\mathrm{Ai}(ze^{2\pi i/3})=0$. By
(9.7.5)--(9.7.6), as $z\to\infty$ uniformly in closed subsectors of $|\mathrm{ph}\,z|<\pi$, with
$\chi:=\frac23z^{3/2}$,
\[
  \mathrm{Ai}(z)\sim\frac{e^{-\chi}}{2\sqrt\pi\,z^{1/4}}\sum_{k\ge0}\frac{(-1)^ku_k}{\chi^k},\qquad
  \mathrm{Ai}'(z)\sim-\frac{z^{1/4}e^{-\chi}}{2\sqrt\pi}\sum_{k\ge0}\frac{(-1)^kv_k}{\chi^k},
\]
where, by (9.7.2), $u_0=v_0=1$, $u_k=(2k+1)(2k+3)\cdots(6k-1)/(216^kk!)$ and $v_k=\frac{6k+1}{1-6k}u_k$; in
particular $u_1=\frac5{72}$ and $v_1=-\frac7{72}$.

\emph{(c) Modified Bessel functions.} (10.25.2)
$I_\nu(z)=(\frac12z)^\nu\sum_{k\ge0}(\frac14z^2)^k/(k!\,\Gamma(\nu+k+1))$, analytic in $\C\setminus(-\infty,0]$
(\S10.25(ii)); (10.27.4) $K_\nu=\frac12\pi(I_{-\nu}-I_\nu)/\sin(\nu\pi)$ for $\nu\notin\mathbb Z$; (10.31.2)
$K_0(z)=-(\ln(\frac12z)+\gamma_{\mathrm E})I_0(z)+\frac14z^2+\dotsb$, $\gamma_{\mathrm E}$ being Euler's constant and the
dots an entire function of $z^2$; (10.28.2) $\mathcal W\{K_\nu,I_\nu\}(z)=1/z$, which with (6.23) of
\S\ref{sec:tool-bessel} gives $\det\Psi_\alpha=1$.

\emph{(d) Large argument} (\S\S10.17, 10.40, $\nu$ fixed). The coefficients are, by (10.17.1) and (10.17.8),
\begin{gather*}
  a_k(\nu)=\frac{(4\nu^2-1^2)(4\nu^2-3^2)\cdots(4\nu^2-(2k-1)^2)}{k!\,8^k},\\
  b_k(\nu)=\frac{(4\nu^2-1^2)\cdots(4\nu^2-(2k-3)^2)\,(4\nu^2+4k^2-1)}{k!\,8^k},
\end{gather*}
with $a_0(\nu)=b_0(\nu)=1$; for instance $a_1(\nu)=(4\nu^2-1)/8$ and $b_1(\nu)=(4\nu^2+3)/8$. Let
$\mathsf s_\pm(z):=\sum_k(\pm1)^ka_k(\nu)z^{-k}$. The expansions (10.40.1), (10.40.2) and (10.40.5) read
\begin{align*}
  I_\nu(z)&\sim e^z(2\pi z)^{-1/2}\mathsf s_-(z),&&|\mathrm{ph}\,z|<\tfrac12\pi,\\
  K_\nu(z)&\sim(\pi/2z)^{1/2}e^{-z}\mathsf s_+(z),&&|\mathrm{ph}\,z|<\tfrac32\pi,\\
  I_\nu(z)&\sim(2\pi z)^{-1/2}\bigl(e^z\mathsf s_-(z)\pm ie^{\pm\nu\pi i}e^{-z}\mathsf s_+(z)\bigr),
    &&-\tfrac12\pi<\pm\mathrm{ph}\,z<\tfrac32\pi,
\end{align*}
uniformly in closed subsectors; and (10.40.3)--(10.40.4) are the expansions of $I_\nu',K_\nu'$ in the sectors of
(10.40.1)--(10.40.2), with $b_k(\nu)$ in place of $a_k(\nu)$ and an overall minus sign for $K_\nu'$. For the Hankel
functions, with $\varpi:=z-\frac12\nu\pi-\frac14\pi$ (10.17.2) and $\delta>0$ arbitrarily small, (10.17.5)
$H^{(1)}_\nu(z)\sim(2/(\pi z))^{1/2}e^{i\varpi}\sum_ki^ka_k(\nu)z^{-k}$ for $-\pi+\delta\le\mathrm{ph}\,z\le2\pi-\delta$,
and (10.17.6) $H^{(2)}_\nu(z)\sim(2/(\pi z))^{1/2}e^{-i\varpi}\sum_k(-i)^ka_k(\nu)z^{-k}$ for
$-2\pi+\delta\le\mathrm{ph}\,z\le\pi-\delta$; (10.17.11)--(10.17.12) give $H^{(1,2)\prime}_\nu(z)\sim\pm
i(2/(\pi z))^{1/2}e^{\pm i\varpi}\sum_k(\pm i)^kb_k(\nu)z^{-k}$ in the same two sectors. None of these expansions
restricts $|\nu|$.
